\documentclass[a4paper,11pt]{article}
\usepackage{exam-style-842-unit}
\examinehead{842 信息化概论 · 第 4 单元练习册}{练习一\quad 软件工程与开发流程}
\begin{document}

\unitcover{842 信息化概论 · 单元练习册}{第 4 单元\quad 软件技术与开发架构}{练习一\quad 软件工程与开发流程（共 100 分）}{约 70 分钟}{软件的地位与软件开发三阶段；软件危机与软件工程的产生；《人月神话》的 5 个关键点；软件工程的代表性定义；软件过程与瀑布模型；软件开发的主要流程；软件项目管理与 CMM。}

\begin{knowmap}
\textbf{软件的地位：}软件是计算机的灵魂，但软件的开发和运行常受到计算机系统的限制，且对计算机系统有着不同程度的依赖性\\[2pt]
\textbf{软件开发三阶段：}1950—1960 程序设计（个体手工劳动）$\rightarrow$ 1960—1970 软件设计（软件作坊·小组合作）$\rightarrow$ 1970 年至今 软件工程时代（工程化生产）\\[2pt]
\textbf{危机与学科：}软件危机（进度超期·预算超支·质量不达标·项目失败）$\Rightarrow$ 1968 年北约（NATO）计算机科学家会议首次提出“软件危机”与“软件工程”\\[2pt]
\textbf{人月神话（布鲁克斯/IBM360）：}外科手术式团队 · 核心概念少而精 · 必要沟通 · 适度文档 · 适度改进\\[2pt]
\textbf{软件工程：}以系统化、规范化、可量化的过程化方法开发和维护软件；核心问题＝软件过程＋软件管理\\[2pt]
\textbf{软件过程与瀑布模型：}软件过程是建造高质量软件所需完成的任务框架；1970 年罗伊斯提出瀑布模型，划分为需求分析、架构设计、详细设计、编码、测试、维护 6 个基本活动\\[2pt]
\textbf{项目管理：}4P（人员·产品·过程·项目）＋软件度量·计划·风险·质量·过程能力·配置管理；CMM 分初始级、可重复级、已定义级、定量管理级、优化级 5 级
\end{knowmap}

\begin{keyterms}
\textbf{软件工程}：将系统化、严格约束的、可量化的方法应用于软件的开发、运行和维护的工程学科。\\
\textbf{软件危机}：软件开发中进度超期、预算超支、质量难保证、项目大量失败等尖锐矛盾的总称。\\
\textbf{人月神话}：布鲁克斯总结 IBM360 操作系统教训的软件工程经典著作，提出 5 个关键点。\\
\textbf{软件过程}：为建造高质量软件所需完成的任务的框架，即形成软件产品的一系列步骤。\\
\textbf{瀑布模型}：按工序将问题化简、按固定次序逐级下落的软件开发模型，1970 年由罗伊斯提出。\\
\textbf{软件生命周期}：软件从需求分析、设计、编码、测试到维护的全过程。\\
\textbf{里程碑}：软件开发过程中设置的阶段性节点，用于明确每阶段的任务与目标。\\
\textbf{软件项目管理}：为使软件项目按预定成本、进度、质量顺利完成，对人员、产品、过程和项目进行分析与管理的活动。\\
\textbf{4P}：软件项目管理的四个管理对象——人员（People）、产品（Product）、过程（Process）、项目（Project）。\\
\textbf{软件度量}：对软件规模、复杂度、开发进度与质量等进行定量测度的方法与活动。\\
\textbf{风险管理}：识别、评估与应对软件项目开发过程中各种风险的管理活动。\\
\textbf{软件质量保证}：为保证软件产品满足规定质量要求而开展的计划、评审、测试等活动。\\
\textbf{软件配置管理}：对软件配置项（程序、文档、数据等）进行标识、控制、变更与审计的管理活动。\\
\textbf{CMM}：能力成熟度模型，把软件开发视为一个过程并分 5 个等级对组织软件过程能力进行评价。
\end{keyterms}

\newpage
\begin{contrast}{软件开发三阶段（每行依次为：时间·生产方式·主要特征）}
\citem{程序设计阶段}{1950—1960 年}{个体手工劳动}{程序为特定目的编制，通用性差，尚无“软件”概念，多用低级语言}
\citem{软件设计阶段}{1960—1970 年}{小组合作的“软件作坊”}{出现高级语言与结构化方法，软件作为产品被广泛使用，但开发方式仍近似个体化}
\citem{软件工程时代}{1970 年至今}{工程化的生产方式}{第三四代语言、数据库、开发工具与环境、网络、分布式、面向对象技术得到应用}
\end{contrast}

\begin{contrast}{软件工程 / 软件过程 / 软件项目管理（每行依次为：定位·关注对象·核心任务）}
\citem{软件工程}{工程学科与总体方法论}{软件的开发、运行和维护全过程}{以系统化、规范化、可量化方法提高质量、降低成本}
\citem{软件过程}{软件工程实施的工作框架}{形成软件产品的一系列步骤、中间产品、资源与角色}{定义并管理“怎么做”，如瀑布模型规定的阶段次序}
\citem{软件项目管理}{软件工程的组织管理维度}{人员、产品、过程和项目（4P）}{按预定成本、进度、质量完成软件交付}
\end{contrast}

\begin{contrast}{CMM 的 5 个等级（每行依次为：过程特征·管理要点·总体评价）}
\citem{初始级}{软件过程无序甚至混乱，对过程几乎没有定义}{管理是反应式的，成功只取决于个人努力}{过程能力不可预测，属最低等级}
\citem{可重复级}{建立基本的项目管理过程，跟踪费用、进度和功能特性}{制定必要的过程纪律}{能重复早先类似应用项目取得的成功}
\citem{已定义级}{软件管理和工程两方面的过程已文档化、标准化}{所有项目使用经批准、剪裁的标准软件过程}{组织有了统一的软件过程标准}
\citem{定量管理级}{收集软件过程和产品质量的详细度量数据}{对软件过程和产品都有定量的理解与控制}{管理由定性走向定量}
\citem{优化级}{过程的量化反馈和先进的新思想、新技术促进过程不断改进}{持续开展过程改进}{过程能力持续提升，为 CMM 最高等级}
\end{contrast}

\qtypename{一、单项选择题}{10 题，每题 3 分，共 30 分}
\q 1968 年第一次提出“软件危机”和“软件工程”概念的会议是（\quad）
\choicesline{A. ISO 年会}{B. 北约计算机科学家会议}{C. IEEE 年会}{D. 世界电子商务会议}

\q 按教材对软件开发阶段的划分，1960—1970 年属于（\quad）
\choicesline{A. 程序设计阶段}{B. 软件设计阶段}{C. 软件工程时代}{D. 面向对象阶段}

\q 为建造高质量软件所需完成的任务的框架，即形成软件产品的一系列步骤，称为（\quad）
\choicesline{A. 软件工程}{B. 软件过程}{C. 软件项目管理}{D. 软件体系结构}

\q 1970 年提出著名的“瀑布模型”的学者是（\quad）
\choicesline{A. 布鲁克斯}{B. 温斯顿·罗伊斯}{C. 迈克尔·波特}{D. 冯·诺依曼}

\q 瀑布模型把软件生命周期划分为 6 个基本活动，其正确顺序是（\quad）
\choices{A. 需求分析$\rightarrow$架构设计$\rightarrow$详细设计$\rightarrow$编码$\rightarrow$测试$\rightarrow$维护}
{B. 需求分析$\rightarrow$详细设计$\rightarrow$架构设计$\rightarrow$编码$\rightarrow$测试$\rightarrow$维护}
{C. 架构设计$\rightarrow$需求分析$\rightarrow$编码$\rightarrow$详细设计$\rightarrow$测试$\rightarrow$维护}
{D. 需求分析$\rightarrow$编码$\rightarrow$详细设计$\rightarrow$架构设计$\rightarrow$测试$\rightarrow$维护}

\q 布鲁克斯写作《人月神话》所总结的实践来自（\quad）
\choicesline{A. IBM360 操作系统}{B. UNIX 操作系统}{C. Windows 操作系统}{D. 阿波罗登月软件}

\q 软件项目管理要进行分析和管理的 4 个要素（4P）是（\quad）
\choices{A. 人员、产品、过程、项目}{B. 人员、资金、设备、进度}
{C. 需求、设计、编码、测试}{D. 计划、组织、领导、控制}

\q 20 世纪 70 年代中期美国国防部研究发现，软件项目不能按时提交、预算超支、质量达不到要求，其中因管理不善引起而非技术原因引起的项目约占（\quad）
\choicesline{A. 30\%}{B. 50\%}{C. 70\%}{D. 90\%}

\q CMM 中软件过程无序甚至混乱、对过程几乎没有定义、成功只取决于个人努力，该等级是（\quad）
\choicesline{A. 初始级}{B. 可重复级}{C. 已定义级}{D. 优化级}

\q 下列不属于软件项目管理内容的是（\quad）
\choicesline{A. 软件项目计划}{B. 风险管理}{C. 软件配置管理}{D. 计算机硬件采购}

\qtypename{二、多项选择题}{4 题，每题 4 分，共 16 分；每题有两个或两个以上正确项}
\q 《人月神话》对软件开发过程提出的关键点包括（\quad）
\choicesline{A. 外科手术式的团队组织}{B. 核心概念少而精}{C. 必要的沟通手段}{D. 保持适度文档}

\q 瀑布模型的优点包括（\quad）
\choicesline{A. 设里程碑明确阶段目标}{B. 可为每阶段制订计划、预算}{C. 通过阶段评审纳入正轨}{D. 严格计划保证按时交付}

\q 瀑布模型的缺点包括（\quad）
\choicesline{A. 缺乏灵活性}{B. 小错误被逐级放大}{C. 返回上级代价高昂}{D. 规模越大成功机率越低}

\q 关于 CMM，下列说法正确的有（\quad）
\choices{A. 全称是能力成熟度模型}{B. 把软件开发视为一个过程}
{C. 共分 5 个等级，优化级最高}{D. 只用于评价个人程序员的编程能力}

\qtypename{三、判断题}{6 题，每题 2 分，共 12 分；正确的画“√”，错误的画“×”并写出改正后的表述}
\q 软件工程就是编写程序的技术，只关注编码环节。\quad（\quad）\quad 改正：\fixblank

\q 瀑布模型规定各阶段自上而下、相互衔接，具有固定次序。\quad（\quad）\quad 改正：\fixblank

\q “软件危机”和“软件工程”是 1970 年温斯顿·罗伊斯提出瀑布模型时第一次提出的。\quad（\quad）\quad 改正：\fixblank

\q 缩短软件项目进度的有效办法，是不断向进度落后的项目中增加人手。\quad（\quad）\quad 改正：\fixblank

\q 软件项目管理的根本目的是让软件项目在管理者的控制之下，以预定成本按期、按质地完成交付。\quad（\quad）\quad 改正：\fixblank

\q CMM 的优化级是能力成熟度模型的最高等级。\quad（\quad）\quad 改正：\fixblank

\qtypename{四、配对题}{6 对，每对 1 分，共 6 分；把 A 栏与 B 栏对应内容配对}
\q 将下列人物、事件或概念与其对应内容配对，把 B 栏字母填入括号。
\par\vspace{0.2em}
\noindent\begin{tabular}{@{}p{0.44\textwidth}@{\hspace{1.2em}}p{0.44\textwidth}@{}}
（1）1968 年北约会议 \hfill（\quad） & A. 首次提出“软件危机”与“软件工程” \\
（2）1970 年温斯顿·罗伊斯 \hfill（\quad） & B. 提出瀑布模型 \\
（3）佛瑞德·布鲁克斯 \hfill（\quad） & C. 总结 IBM360 教训，著《人月神话》 \\
（4）SW-CMM \hfill（\quad） & D. 能力成熟度模型的 5 个等级 \\
（5）软件过程 \hfill（\quad） & E. 形成软件产品的一系列步骤 \\
（6）软件配置管理 \hfill（\quad） & F. 软件项目管理的内容之一 \\
\end{tabular}\par\vspace{0.3em}

\qtypename{五、简答题}{3 题，每题 8 分，共 24 分}
\q 什么是软件工程？写出教材中两种有代表性的定义，并说明软件工程要解决的两个重要问题。\anslines{2.4cm}

\q 什么是软件危机？简述软件危机的主要表现与产生原因，并说明软件工程产生的背景。\anslines{2.4cm}

\q 简述瀑布模型的核心思想与 6 个基本活动，并分别说明其优点和缺点。\anslines{2.4cm}

\newpage
\qtypename{六、综合题}{1 题，共 12 分}
\q 某高校团队为县农业局开发“农产品质量追溯管理系统”，原计划 6 个月完成，实际拖到 15 个月；项目经理不断增派程序员，进度反而更慢；系统临近上线又因需求频繁变更而大量返工。请运用《人月神话》的 5 个关键点、瀑布模型的特点以及软件项目管理的内容，分析该项目陷入困境的原因，并提出至少 4 条改进建议。\anslines{5cm}

\end{document}
